\documentclass[10pt,a4]{examenEPN0}
\usepackage{comandosEPN}

\usepackage{empheq}

\newcommand{\itemcp}{\stepcounter{enumi}\item[ \textcolor{red}{\theenumi.$^*$}] }

\newcommand{\iitemcp}{\stepcounter{enumii}\item[ \textcolor{red}{\theenumii)$^*$}] }


\materia{Álgebra Lineal}
%\autor{Departamento de Formación Básica}
\autor{GR \hspace{1cm}.}
%\fecha{Semestre 2026-A}
\fecha{Nombre: \hspace{9cm} Cédula:}
\logouno[1.6cm]{Logo_EPN.eps}
\logodos[2.6cm]{Logo_DFB.eps}

\newcommand{\la}{\lambda}

%Datos
\newcommand{\datosAlgebra}[2]{
\EnExamen
{
\vspace{5mm}
    \noindent
    Nombres: \hrulefill\
    Carrera: \rule{4.5cm}{0.4pt}\
    \ifthenelse{\equal{#1}{}}%
    {Paralelo: \underline{\parbox{1.8cm}{\parbox{1cm}{}}}\\[3mm]}
    {Paralelo: \underline{\parbox{1.8cm}{\centering #1}}\\[3mm]}
    C\'edula: \rule{5cm}{0.4pt}\
    Firma: \hrulefill\
    Fecha: \underline{\parbox{2.1cm}{\centering #2}}
}}


% \newcommand{\indicacionesPrueba}{%
% \EnExamen{
% \begin{indicaciones}
%     \item La prueba tiene una duraci\'on  de \textbf{120 minutos.}
%     \item No se permite la utilizaci\'on de dispositivos electr\'onicos (celular, tablets, calculadora, etc.)
%     \item Leer detenidamente los enunciados, \textbf{justificar la respuesta para acceder al puntaje total}.
%     \item La prueba es completamente individual. Los estudiantes deberán dejar sus teléfonos celulares guardados y no podrán tenerlos consigo (por ejemplo, en los bolsillos de su ropa) durante toda la evaluación. En caso de detectarse algún acto de deshonestidad académica o el incumplimiento de esta norma, no podrán seguir rindiendo la prueba y su calificación será \textbf{cero}.
% \end{indicaciones}
% \pagestyle{empty}
% }}

\newcommand{\indicacionesPrueba}{%
\EnExamen{
\begin{indicaciones}
    \item La prueba tiene una duraci\'on  de \textbf{60 minutos.}
    \item Leer detenidamente los enunciados, \textbf{justificar la respuesta} para acceder al puntaje total.
    \item No se permite la utilizaci\'on de dispositivos electr\'onicos (celular, tablets, calculadora, etc.)
    \item La prueba es completamente individual
    \item No se permite portar teléfonos celulares, relojes inteligentes ni otros dispositivos electrónicos durante el examen, aun cuando permanezcan en los bolsillos. Estos dispositivos deberán guardarse en la mochila antes del inicio de la evaluación y esta deberá mantenerse cerrada durante todo el examen.
    \item He leído las indicaciones y acepto que el incumplimiento de estas se penaliza con \textbf{cero}. Firma:...................................

    
\end{indicaciones}
\pagestyle{empty}
}}

\newcommand{\indicacionesExamen}{
\EnExamen{\begin{indicaciones}
    \item La evaluación tiene una duraci\'on  de \textbf{120 minutos.}
    \item No se permite la utilizaci\'on de dispositivos electr\'onicos (celular, tablets, calculadora, etc.)
    %\item Leer detenidamente los enunciados, justificar la respuesta para acceder al puntaje total.
    \item El examen es completamente individual. En caso de detectarse alg\'un acto de deshonestidad acad\'emica no podr\'a seguir rindiendo la prueba y su calificaci\'on ser\'a \textbf{cero}.
\end{indicaciones}
}\pagestyle{empty}
}

%% Imprimir las una hoja A4, dos veces en A5
\newcommand{\imprimirDoble}
{%
\usepackage{pgfmorepages}
\pgfmorepagesloadextralayouts
\pgfpagesuselayout{repeated 2-up}[a4paper,landscape]
\thispagestyle{empty}
}


%Este sí se puede comentar
\newenvironment{solucion}{\color{blue}\proof[Solución]}{\endproof\color{black}}




% \renewcommand{\@autor}{
% ---\\
% \phantom{.} 
% ...}
% \makeatother

\examen{Examen 01: solucionario y rúbrica}
\makeatletter
\renewcommand{\@autor}{Departamento de Formación Básica}
\renewcommand{\@fecha}{GR6M -- Viernes de 11 a 13}
\makeatother
\begin{document}
\encabezado
\noindent\textbf{Criterio general (20 puntos).} Se asigna crédito parcial por
los pasos correctos. Un error aritmético se penaliza donde se origina;
los pasos posteriores se valoran por su coherencia, sin repetir la penalización.
Se aceptan procedimientos equivalentes debidamente justificados.
\par\noindent\textbf{Distribución de puntajes.} El enunciado no especifica
puntajes; esta rúbrica asigna 4 puntos a cada pregunta.
La fecha del archivo fuente es 17/11/2026; se conserva sin modificar el examen.

\section*{1. Sistema con parámetro (4 puntos)}
\noindent\textbf{Rúbrica.}
\begin{itemize}[itemsep=2pt]
    \item Calcular el determinante con procedimiento: 1,5 puntos; establecer cuándo hay solución única: 0,5 puntos.
    \item Sustituir $\alpha=3$ y reducir: 1 punto; exhibir una contradicción: 0,5 puntos; concluir que el conjunto solución es vacío: 0,5 puntos.
\end{itemize}
Mediante $F_2\leftarrow F_2-2F_1$ y $F_3\leftarrow F_3-F_1$,
\[
    \det(A)=\det\begin{pmatrix}1&-\alpha&4\\0&\alpha-3&2\alpha+1\\0&0&\alpha+5\end{pmatrix}
    =(\alpha-3)(\alpha+5).
\]
Hay solución única exactamente para $\alpha\notin\{3,-5\}$.
Para $\alpha=3$, las mismas operaciones sobre las ecuaciones producen
$7x_3=2$ y $8x_3=-1$. Como $2/7\ne-1/8$, el sistema es
incompatible y su conjunto solución es $\boxed{\varnothing}$.

\section*{2. Combinaciones lineales y espacio generado (4 puntos)}
\noindent\textbf{Rúbrica.}
\begin{itemize}[itemsep=2pt]
    \item Plantear la combinación y hallar los coeficientes: 0,5 puntos; obtener $k$: 0,5 puntos.
    \item Describir el espacio generado: 0,5 puntos; justificar que es un plano: 0,5 puntos.
    \item Comprobar la pertenencia del vector dado: 0,75 puntos; concluir: 0,25 puntos.
    \item Justificar la igualdad de los espacios generados para el $k$ obtenido: 1 punto.
\end{itemize}
Si $v_3=a v_1+b v_2$, las primeras coordenadas dan
$a+2b=1$ y $2a+b=-4$. Por tanto, $a=-3$, $b=2$ y
$k=a-b=-5$; en ese caso $v_3=-3v_1+2v_2$.
Como $v_1\times v_2=(-3,3,-3)\ne0$,
\[
    \operatorname{span}\{v_1,v_2\}
    =\{a(1,2,1)+b(2,1,-1):a,b\in\R\}
    =\{(x,y,z):x-y+z=0\}.
\]
Es un plano por el origen. Para $v=(4,5,1)$ se tiene
$4-5+1=0$, por lo que \textbf{sí pertenece}; explícitamente,
$v=2v_1+v_2$.
Para $k=-5$, añadir $v_3$, que ya es combinación de $v_1,v_2$,
no cambia el espacio generado. Una inclusión se sigue de esa combinación
y la otra de que $v_1,v_2$ están entre los tres generadores.

\section*{3. Intersección de subespacios de polinomios (4 puntos)}
\noindent\textbf{Rúbrica.}
\begin{itemize}[itemsep=2pt]
    \item Resolver las condiciones sobre $a,b,c$: 1 punto; expresar la intersección: 1 punto.
    \item Plantear las condiciones sobre $k$: 1 punto; hallar el valor y verificar: 1 punto.
\end{itemize}
Las condiciones $a+b=0$, $b+c=0$ y $a-c=0$ equivalen a
$a=c=t$ y $b=-t$. Así,
\[
    \boxed{S\cap T=\{t(x^2-x+1):t\in\R\}.}
\]
Para el polinomio dado, $a+b=k-2$, $b+c=2k-4$ y
$a-c=2-k$. Se anulan simultáneamente para $k=2$, y el
polinomio es $3x^2-3x+3$.

\section*{4. Matrices de traza cero (4 puntos)}
\noindent\textbf{Rúbrica.}
\begin{itemize}[itemsep=2pt]
    \item Verificar que la matriz cero pertenece a $W$: 1 punto.
    \item Demostrar cerradura bajo suma por linealidad de la traza: 1,5 puntos.
    \item Demostrar cerradura bajo producto por escalar: 1 punto.
    \item Concluir que es subespacio: 0,5 puntos.
\end{itemize}
La matriz cero tiene traza cero. Si $A,B\in W$ y $\lambda\in\R$,
\[
    \operatorname{tr}(A+B)=\operatorname{tr}(A)+\operatorname{tr}(B)=0,
    \qquad \operatorname{tr}(\lambda A)=\lambda\operatorname{tr}(A)=0.
\]
Por tanto, $W$ es subespacio de $\R^{3\times3}$. También se acepta
identificar $W$ como núcleo de la aplicación lineal traza.

\section*{5. Servicios del taller (4 puntos)}
\noindent\textbf{Rúbrica.}
\begin{itemize}[itemsep=2pt]
    \item Definir las variables y unidades: 0,25 puntos; plantear las ecuaciones: 0,75 puntos (0,25 por ecuación); identificar $A$, $\mathbf{x}$ y $b$: 1 punto.
    \item Resolver con procedimiento: 1,5 puntos; interpretar y comprobar las tarifas: 0,5 puntos.
\end{itemize}
Sean $d,a,l$ las tarifas de diagnóstico, ajuste y limpieza, en dólares.
\[
    d+a+l=60,\qquad a+l=35,\qquad d-l=10.
\]
\[
    A=\begin{pmatrix}1&1&1\\0&1&1\\1&0&-1\end{pmatrix},\quad
    \mathbf{x}=\begin{pmatrix}d\\a\\l\end{pmatrix},\quad
    b=\begin{pmatrix}60\\35\\10\end{pmatrix}.
\]
Restando las dos primeras ecuaciones, $d=25$. De $d-l=10$,
$l=15$; luego $a=20$. Las tarifas son \textbf{\$25 por diagnóstico,
\$20 por ajuste y \$15 por limpieza}. Se comprueba
$25+20+15=60$, $20+15=35$ y $25-15=10$.
\end{document}
